% Application : inférence, disponibilité, live, API, interface, traçabilité (partie 5 ; ADR-0040 à 0046).

\chapter{Application}
\label{chap:application}

\section{Inférence : la même fonction qu'à l'entraînement}

L'inférence n'a aucune formule de variable propre : pour prédire les matchs d'un jour \(J\), elle appelle
la fonction qui construit le jeu d'entraînement sur l'historique arrêté à \(J\), un jour à la fois
(ADR-0040). Sur 17 jours de contrôle, les 622 lignes produites ainsi sont identiques au bit près à celles
du jeu d'entraînement. Deux modes existent :

\begin{itemize}
  \item \textbf{rejeu} : un match passé, prédit par le modèle du pli de sa saison, qui ne l'a jamais vu ;
        les prédictions reproduisent celles de l'évaluation à la précision machine. Les saisons 2021-22 à
        2024-25 sont ouvertes ; les saisons scellées restent fermées jusqu'au test de l'horizon H2
        (ADR-0045) ;
  \item \textbf{live} : un match à venir, prédit par le modèle final, et seulement avant son coup
        d'envoi.
\end{itemize}

\section{Disponibilité}

Pour chaque variable requise, l'inférence dit si elle est \emph{présente}, \emph{manquante} (avec la
raison) ou \emph{périmée} (source en retard). Une seule variable non présente suffit : le match est
\code{unavailable}, et aucune prédiction n'est faite. Aucune valeur de remplacement n'est jamais
utilisée, car une moyenne mise à la place d'une variable absente donnerait une prédiction fausse
présentée avec la même assurance qu'une vraie.

\section{Le live sans abonnement}

Après le gel, le calendrier d'API-FOOTBALL reste figé (dates souvent fausses, aucun résultat). Les
résultats et les prochains matchs de football-data lui sont \emph{superposés en mémoire}, sans rien
écrire dans le référentiel (ADR-0042). L'appariement se fait par championnat, saison et équipes, sans la
date ; un nom inconnu ou une clé ambiguë n'est jamais deviné. Les tirs de football-data remplacent ceux
d'API-FOOTBALL : l'écart de log-loss mesuré est négligeable (\(-0{,}00002\), \([-0{,}00021 ; 0{,}00018]\),
ADR-0041).

\paragraph{Fraîcheur.} football-data publie les résultats environ deux fois par semaine. Avec la règle
d'origine (un championnat entièrement à jour), 91~\% des matchs des sept jours suivants étaient
indisponibles lors du premier live, le 22 octobre. La règle a été révisée après mesure (ADR-0046) : un
match n'attend plus que les résultats de ses deux équipes. L'Elo reste exact ; les moyennes glissantes
deviennent approchées, à moins de 0,0011 près sur \(\lambda\) d'après une répétition sur 2023-24. La part
d'indisponibles est tombée à 14~\%, puis à 10~\% le 31 octobre.

\section{API et interface}

L'API FastAPI (7 routes) liste les matchs d'une date avec leur disponibilité, calcule et trace une
prédiction, décrit le modèle actif et la fraîcheur des données. Elle n'écoute que sur \code{127.0.0.1}
et refuse toute date scellée (code 403). L'interface Streamlit a trois écrans :

\begin{itemize}
  \item \textbf{Matchs} : mode (rejeu ou live), date, championnats ; une pastille de disponibilité par
        match, avec sa raison ; un bouton \og Prédiction \fg{} actif seulement si le match est
        disponible ;
  \item \textbf{Détail} : loi du total en barres, \(E[T]\), \(\lambda\) de chaque équipe, \(P(T > 2{,}5)\),
        intervalle et couverture annoncée, variables utilisées, version du modèle et des données, score
        réel en rejeu, référence de marché en live ;
  \item \textbf{Modèle} : carte d'identité, métriques de validation, courbe d'apport, fraîcheur.
\end{itemize}

L'interface ne parle qu'à l'API, en HTTP : un test d'architecture lui interdit d'importer la logique
métier ou la base.

\section{Traçabilité et journal prospectif}

Chaque prédiction créée est écrite dans \code{ops.prediction} avec tout ce qui a servi à la faire :
modèle, version des données, variables, loi, intervalle, date de complétude des sources. L'écriture est
idempotente : une seconde demande renvoie la ligne existante. En live, une prédiction n'est acceptée
qu'avant le coup d'envoi, et une réponse indisponible n'est jamais écrite. Au 31 octobre 2026, 89
prédictions live sont tracées. Ce journal prospectif servira à évaluer le modèle sur des matchs
réellement futurs ; il n'est pas encore évalué, car lire leurs scores reviendrait à ouvrir le scellé au
développement (ADR-0045).

\section{Tests de bout en bout}

Un test automatisé, en intégration continue, parcourt toute la chaîne sur une petite ligue synthétique :
brut, \code{load}, porte des données, inférence, API, interface, ligne dans \code{ops.prediction}. Sur
les vraies données, la même chaîne a été vérifiée en rejeu et en live le 31 octobre 2026 : l'écran
\og Matchs \fg{} s'affiche en 12 à 17 secondes à froid et en moins d'une demi-seconde ensuite, une
prédiction en 2 à 5 secondes.
